\section{Marco teórico y trabajos relacionados}

La transformación digital en la agricultura colombiana se enmarca en un contexto global de adopción de tecnologías emergentes. Barrios-Ulloa et al. (2025) analizan cómo las redes 6G pueden modernizar la agricultura en Colombia, destacando oportunidades en IoT, IA y conectividad rural. Mohammad Raziuddin Chowdhury et al. (2023) exploran el rol de la agricultura digital en la conversión de áreas rurales en "smart villages", enfatizando el desarrollo sostenible.

Eugênio Pacceli Reis da Fonseca et al. (2019) presentan un sistema de información verde para la gestión sostenible de agroecosistemas, útil para integrar datos ambientales. Amit Degada et al. (2021) examinan la transformación digital en aldeas inteligentes mediante IoT, abordando el uso de sensores y automatización en contextos rurales.

La Comisión Económica para América Latina y el Caribe (CEPAL, 2020) proporciona un análisis general del impacto de las TIC en la agricultura y el desarrollo rural. Investigaciones sobre "Retos de la Agricultura Digital" (Revista Colombiana de Investigaciones Agroindustriales, 2025) estudian barreras de digitalización en Colombia, incluyendo infraestructura y desigualdades.

Estudios como el de "How do coffee farmers engage with digital technologies?" (2024) analizan la adopción tecnológica en caficultores colombianos desde una perspectiva de capacidades. El artículo sobre "Digital integration toward smart and sustainable agriculture" (2025) examina impactos de la digitalización en la sostenibilidad agrícola europea.

Artículos sobre blockchain en agricultura sostenible en Sudamérica (2024) y mercados campesinos en Cundinamarca (2022) ofrecen insights sobre comercio justo y circuitos cortos de comercialización. Finalmente, documentos del IICA (2025) y FAO (2023) respaldan la importancia de la digitalización en la asistencia técnica y el fortalecimiento rural.

Estos trabajos identifican vacíos en la aplicación contextual de tecnologías digitales en el agro colombiano, especialmente en la eliminación de brechas digitales y la promoción de modelos inclusivos.